\documentclass[letterpaper,11pt]{article}

% Shared LaTeX preamble for CV and Cover Letter

\usepackage{iftex}

\ifPDFTeX
  \usepackage[T2A]{fontenc}
  \usepackage[utf8]{inputenc}
\else
  \usepackage{fontspec}
\fi

\usepackage[english, russian]{babel}

\usepackage{latexsym}
\usepackage[empty]{fullpage}
\usepackage{titlesec}
\usepackage{marvosym}
\usepackage[usenames,dvipsnames]{color}
\usepackage{verbatim}
\usepackage{enumitem}
\usepackage[hidelinks]{hyperref}
\usepackage{fancyhdr}
\usepackage{tabularx}
\usepackage{graphicx}

\input{glyphtounicode}
\pdfgentounicode=1

\pagestyle{fancy}
\fancyhf{} % clear all header and footer fields
\fancyfoot{}
\renewcommand{\headrulewidth}{0pt}
\renewcommand{\footrulewidth}{0pt}

\urlstyle{same}
\raggedbottom
\raggedright

\newcommand{\username}{alexo375}
\newcommand{\github}{https://github.com/\username}
\newcommand{\cpp}{C\texttt{++}}


% Adjust margins
\addtolength{\oddsidemargin}{-0.5in}
\addtolength{\evensidemargin}{-0.5in}
\addtolength{\textwidth}{1in}
\addtolength{\topmargin}{-.7in} % было "-.5in". Подняли текст чуть выше
\addtolength{\textheight}{1.4in} % было "1.0in". Увеличили рабочую область страницы

\setlength{\tabcolsep}{0in}

% Sections formatting
\titleformat{\section}{
  \vspace{-4pt}\scshape\raggedright\large
}{}{0em}{}[\color{black}\titlerule \vspace{-5pt}]

%-----------CUSTOM-COMMANDS-----------
\newcommand{\resumeItem}[1]{
  \item\small{
    {#1 \vspace{-2pt}}
  }
}

\newcommand{\resumeSubheading}[4]{
  \vspace{-2pt}\item
    \begin{tabular*}{0.97\textwidth}[t]{l@{\extracolsep{\fill}}r}
      \textbf{#1} & #2 \\
      \textit{\small#3} & \textit{\small #4} \\
    \end{tabular*}\vspace{-7pt}
}

\newcommand{\resumeSubheadingCompany}[2]{
  \vspace{-2pt}\item
    \begin{tabular*}{0.97\textwidth}[t]{l@{\extracolsep{\fill}}r}
      \textbf{#1} & #2 \\
    \end{tabular*}\vspace{-2pt}
}

\newcommand{\resumeSubrole}[1]{
  \vspace{4pt}
  \textit{\small#1}\par\vspace{2pt}
}

\newcommand{\resumeProjectHeading}[2]{
    \item
    \begin{tabular*}{0.97\textwidth}{l@{\extracolsep{\fill}}r}
      \small#1 & #2 \\
    \end{tabular*}\vspace{-7pt}
}

\newcommand{\resumeSubHeadingListStart}{\begin{itemize}[leftmargin=0.15in, label={}]}
\newcommand{\resumeSubHeadingListEnd}{\end{itemize}}
\newcommand{\resumeItemListStart}{\begin{itemize}}
\newcommand{\resumeItemListEnd}{\end{itemize}\vspace{-5pt}}

\begin{document}

%----------ЗАГОЛОВОК----------
\begin{center}
  \begin{minipage}[b]{0.80\textwidth}
      \textbf{\Huge \scshape Горобченко Александр} \\ \vspace{4pt}
      \large ML Engineer \\ \vspace{4pt}
      \small \href{tel:+79953755007}{+7 (995) 375-50-07} $|$ 
      \href{mailto:\username@yandex.ru}{\underline{\username@yandex.ru}} $|$ 
      \href{https://t.me/\username}{\underline{t.me/\username}} $|$
      \href{\github}{\underline{github.com/\username}} \\ \vspace{2pt}
      \small 11.09.2004 \\ \vspace{2pt}
      \small Рязань / Готов к релокации в Санкт-Петербург
  \end{minipage}%
  \begin{minipage}[b]{0.18\textwidth}
      \raggedleft
      \includegraphics[width=4cm, keepaspectratio]{me.png} 
  \end{minipage}
\end{center}

%----------О-СЕБЕ----------
\section{О Себе}
 \begin{itemize}[leftmargin=0.15in, label={}]
    \small{\item{
     ML Engineer с экспертизой в Computer Vision и оптимизации алгоритмов (Python / \cpp).
     Имею 3 года R\&D-опыта в области обработки изображений и научные публикации.
     Специализируюсь на полном цикле создания ML-решений: от исследований и подготовки датасетов до оптимизации 
     инференса и микросервисного деплоя.
    }} 
 \end{itemize}

%-----------ОПЫТ-----------
\section{Опыт работы}
  \resumeSubHeadingListStart
    \resumeSubheadingCompany
      {\href{https://rsreu.ru/vuz/structure/nii-foton}{\underline{НИИ обработки аэрокосмических изображений <<Фотон>>}}}{Июль 2023 -- Июль 2026}
      
      \resumeSubrole{ML-инженер}
      \small{Отвечал за полный цикл создания ML-компонента системы обнаружения изменений на аэрокосмических снимках площадной съемки: 
      сформировал и разметил обучающий датасет по данным космических аппаратов, обучил нейросетевую модель сегментации
      UNet\plusplus~на ResNet-34 с помощью PyTorch. 
      Разработал прикладное ПО на \cpp~(с GUI на Qt), интегрировав в него обученную модель в формате ONNX и
      объединив детерминированные алгоритмы компьютерного зрения с методами глубокого обучения.}
      
      \vspace{4pt}

      \resumeSubrole{R\&D-инженер}
      \small{Проводил научно-исследовательские работы в области обработки аэрокосмических изображений.
      Устранил появление артефактов при сшивке мозаичных изображений сложных сцен,
      разработав алгоритм формирования стоимостной карты штрафов для корректного проведения линии пореза.
      Выполнил анализ предметной области, по материалам исследований подготовил и опубликовал 4 научные статьи.}
  \resumeSubHeadingListEnd

%-----------ПРОЕКТЫ-----------
\section{Проекты}
    \resumeSubHeadingListStart
      \resumeProjectHeading
          {\href{\github/change-detector}{\textbf{\underline{ПО автоматического обнаружения изменений на спутниковых снимках}}} $|$ \emph{Python, PyTorch, \cpp}}{2026~г.}
          \resumeItemListStart
            \resumeItem{Обучил CNN (UNet\plusplus, ResNet-34) для сегментации изменений на снимках. Метрики: $F1 \approx 0.86$, $IoU \approx 0.76$.}
            \resumeItem{Выполнил ручную разметку кастомного датасета, получив более 2000 эталонных снимков.}
            \resumeItem{Разработал ПО с GUI на \cpp~и QT, в который экспортировал обученную модель в формате ONNX.}
          \resumeItemListEnd
          
      \resumeProjectHeading
          {\href{\github/acc-2025-validation-pipeline}{\textbf{\underline{Веб приложение валидации данных на основе ML}}} $|$ \emph{Python, FastAPI, ML Pipeline}}{2025~г.}
          \resumeItemListStart
            \resumeItem{Реализовал микросервисный пайплайн для ML-валидации изображений в команде из 3 человек.}
            \resumeItem{Интегрировал ML-модели (MediaPipe, DeepFace, Facer, U2Net) для валидации и комплексной оценки изображений по множеству контекстных параметров.}
            \resumeItem{Создал на FastAPI микросервис, обеспечивающий стабильную проверку качества данных.}
          \resumeItemListEnd

      \resumeProjectHeading
          {\href{\github/native-compiler}{\textbf{\underline{Консольная утилита-компилятор}}} $|$ \emph{\cpp}}{2025~г.}
          \resumeItemListStart
            \resumeItem{Разработал CLI-утилиту, демонстрирующую работу основных фаз компилятора с реализацией 
            лексического, семантического, синтаксического анализа; генератора 3-адресного кода и постфиксной формы.}
          \resumeItemListEnd
    \resumeSubHeadingListEnd

%-----------ОБРАЗОВАНИЕ-----------
\section{Образование}
  \resumeSubHeadingListStart
    \resumeSubheading
      {Центральный Университет}{Москва}
      {Математика и компьютерные науки, трек ML $|$ Магистратура}{2025 - н.в.}
    \resumeSubheading
      {Рязанский Государственный Радиотехнический Университет}{Рязань}
      {Факультет вычислительной техники, Программная инженерия $|$ Бакалавриат}{2022 - 2026}
  \resumeSubHeadingListEnd

%-----------НАВЫКИ-----------
\section{Навыки}
 \begin{itemize}[leftmargin=0.15in, label={}]
    \small{\item{
     \textbf{Python, ML \& CV}{: PyTorch, CatBoost, XGBoost, LightGBM, Scikit-Learn, OpenCV, Pandas, NumPy, SciPy} \\
     \textbf{\cpp}{: CMake, ONNX, Boost, GTest, Valgrind, Conan} \\
     \textbf{Data Engineering}{: PostgreSQL, Greenplum, ClickHouse, Amazon S3, Airflow, Spark, Kafka} \\
     \textbf{DevOps \& CI/CD}{: Docker, Git, Github Actions, Linux (NixOS, Ubuntu)} \\
     \textbf{Языки}{: Английский (B2)}
    }}
 \end{itemize}

%-------------------------------------------
\end{document}
